%%%%%%%%%%%%%%%%%%%%%%%%
% $Autor: MansourAhmadyPhoulady, RanjeetsingTawar, HemanthPrabhakaran$
% $Pfad: Car_Licence_Plate_Identification_with_a_JetsonNano$
% $Version: 4.0.0 $
%
% !TeX encoding = utf8
%
%%%%%%%%%%%%%%%%%%%%%%%%

\chapter{Data Version Control}

Effective management of data, models, and experiments is critical in data science, machine learning, and collaborative software development projects. Data Version Control (DVC) extends the capabilities of traditional version control systems, enabling the seamless tracking of large data files and machine learning artifacts. By integrating DVC into a project workflow, teams can ensure reproducibility, facilitate collaboration, and optimize resource management.

\section{Why Data Version Control?}

Implementing a data version control strategy offers several key benefits for machine learning projects:

\begin{enumerate}
	\item \textbf{Tracking Changes:} Datasets and source code evolve throughout the lifecycle of a project. DVC allows you to monitor these changes, providing the ability to revert to earlier versions if necessary and ensuring transparency in how data transforms over time.
	\item \textbf{Facilitating Collaboration:} In team projects, consistent use of datasets and models across all members is essential. DVC helps synchronize data and code, reducing conflicts and ensuring a unified project workflow.
	\item \textbf{Streamlining Experimentation:} Machine learning requires frequent experimentation. By linking data versions to specific code changes, DVC ensures reproducibility and simplifies the comparison of experimental results.
\end{enumerate}

\section{Implementation Plan}

The implementation of data version control in this project involved a structured approach to ensure consistency and collaboration:

\subsection*{Version Control for Project Files}
The first step was establishing GitHub as the central version control system for all project-related files, including the report, manual, and presentations. Team members were required to install Git and align their configurations to match the project's GitHub repository version. The repository served as the primary source for updates, and a communication protocol was established: team members would notify the group before making updates to prevent conflicts during simultaneous fetch and push operations.

\subsection*{Software Version Alignment}
To maintain consistency across computational environments, the team standardized software versions for all tools and packages used in the project. Aligning versions ensured that results could be replicated without discrepancies caused by variations in software behavior. A summary of aligned versions is provided in Table \ref{tab:software-versions}.

\subsection*{Dataset Version Control}
To avoid data corruption and conflicts, dataset version control was centralized on a single team member's system. Weekly backups were created, preserving previous versions to safeguard against accidental overwrites or errors in new versions. This approach ensured that the dataset and trained models remained accessible and reproducible throughout the project.

\begin{table}[h!]
	\centering
	\caption{Software and Package Versions}
	\label{tab:software-versions}
	\begin{tabular}{|l|l|}
		\hline
		\textbf{Software/Package} & \textbf{Version} \\
		\hline
		GitHub                   & Version 3.4.12 (x64) \\
		TeXstudio                & TeXstudio 4.7.3     \\
		PyCharm                  & 2024.2.2            \\
		Python                   & 3.13  (x64)             \\
		TensorFlow               & 2.16.1              \\
		NumPy                    & 1.25.2              \\
		pandas                   & 2.2.3               \\
		matplotlib.pyplot        & 3.10.0               \\
		os (Built-in Python)     & 3.12.1             \\
		\hline
	\end{tabular}
\end{table}

\section{Installation Steps}

To implement DVC in the project workflow, the following steps were followed:

\begin{enumerate}
	\item \textbf{Install DVC:} DVC was installed using \texttt{pip}, \texttt{conda}, or system-specific methods. Detailed installation instructions are available in the official DVC documentation: \url{https://dvc.org/doc/install}.
	
	\item \textbf{Set Up GitHub Repository:} A dedicated GitHub repository was created to house project files, including DVC metadata. Team members were invited to the repository and granted appropriate access permissions.
	
	\item \textbf{Initialize DVC:} Within the project directory, DVC was initialized using the \texttt{dvc init} command. This step created the necessary configuration files for managing data and metadata.
	
	\item \textbf{Add Data to DVC:} Data files were added to DVC using the \texttt{dvc add} command, which generated corresponding metafiles. These files were subsequently tracked using Git, ensuring synchronization between DVC and GitHub.
	
	\item \textbf{Configure Remote Storage:} A remote storage solution (e.g., Amazon S3, Google Cloud Storage) was configured to store large data files and enable easy sharing. The configuration process followed guidelines provided in the DVC documentation: \url{https://dvc.org/doc/user-guide/setup-remote-storage}.
	
	\item \textbf{Commit and Push Changes:} Changes to the project directory, including DVC metadata, were committed using Git commands (\texttt{git add}, \texttt{git commit}, \texttt{git push}) to ensure synchronization across the team.
	
	\item \textbf{Collaborate:} Team members cloned the repository and used DVC commands to pull datasets and track updates. Regular communication ensured smooth collaboration and minimized conflicts.
\end{enumerate}

By adopting these practices, the project successfully managed datasets, models, and experiments, ensuring reproducibility and streamlined collaboration across the team.
